\documentclass[12pt]{article}
\usepackage[T1]{fontenc}
\usepackage[utf8]{inputenc}
\usepackage{lmodern}
\usepackage{textcomp}
\usepackage{enumitem}
\usepackage{geometry}
\geometry{margin=1in}
\begin{document}
\begin{center}
Answer Key - Science - K-12 Level Assessment 12 \
\small (Answers are ordered exactly as in the question paper.)
\end{center}

\section*{Multiple Choice}
\begin{enumerate}[label=\arabic*.]
\item \textbf{Answer:} B
\\ \textit{Options:} A) Drop each object in different liquids.; B) Record in a table how each object behaves.; C) Measure the length and mass of each object.; D) Test the objects in groups from smallest to largest.
\\ \textit{Note:} Recording the results in a table allows Jake and Molly to systematically compare the behaviors of their objects in a clear and organized manner, facilitating easier analysis of which objects can float.
\item \textbf{Answer:} B
\\ \textit{Options:} A) Meltwater decreases, and the size of the glacier decreases.; B) Meltwater decreases, and the size of the glacier increases.; C) Meltwater increases, and the size of the glacier decreases.; D) Meltwater increases, and the size of the glacier increases.
\\ \textit{Note:} When winter and summer temperatures are below normal and precipitation rates are above normal, increased snowfall accumulates on glaciers, leading to a decrease in meltwater and an overall increase in the size of the glacier.
\item \textbf{Answer:} D
\\ \textit{Options:} A) root; B) cone; C) trunk; D) needle
\\ \textit{Note:} The needles of a pine tree contain chlorophyll and perform photosynthesis, allowing them to convert sunlight into food for the tree.
\item \textbf{Answer:} C
\\ \textit{Options:} A) The Sun is a very large star that exists far from the Milky Way Galaxy.; B) The Sun is at the center of the Milky Way Galaxy.; C) The Sun is a medium-sized star near the edge of the Milky Way Galaxy.; D) The Sun has the Milky Way and several other galaxies moving in orbits around it.
\\ \textit{Note:} The Sun is classified as a G-type main-sequence star (or yellow dwarf) and is located about 26,000 light-years from the center of the Milky Way Galaxy, making it a medium-sized star in terms of both size and position within the galaxy.
\item \textbf{Answer:} B
\\ \textit{Options:} A) the presence of a fluid mantle; B) the presence of Van Allen belts; C) the presence of gravitational force; D) the presence of water in three states
\\ \textit{Note:} The presence of Van Allen belts indicates that a planet has a magnetic field, which is typically generated by the movement of molten iron in its core, suggesting that the planet likely has an iron core similar to Earth's.
\end{enumerate}

\section*{True / False}
\begin{enumerate}[label=\arabic*.]
\setcounter{enumi}{5}
\item \textbf{Answer:} True
\\ \textit{Options:} A) True; B) False
\\ \textit{Note:} Saccharides, which are simple sugars, chemically bond together through glycosidic linkages to form polysaccharides, thereby categorizing them as complex carbohydrates due to their larger molecular structure and varied functions in biological systems.
\item \textbf{Answer:} True
\\ \textit{Options:} A) True; B) False
\\ \textit{Note:} The statement is true because in a solid, the particles are tightly packed together and can only vibrate in place, maintaining a fixed structure similar to spectators who remain seated in assigned spots.
\item \textbf{Answer:} False
\\ \textit{Options:} A) True; B) False
\\ \textit{Note:} The momentum of an object is calculated by multiplying its mass by its velocity, so for a 0.15 kg baseball moving at 40 m/s, the momentum is 6 kg·m/s, not 0.025 kg·m/s.
\item \textbf{Answer:} True
\\ \textit{Options:} A) True; B) False
\\ \textit{Note:} A population refers to a group of individuals of the same species living in a specific area, whereas a community encompasses all the different species interacting in that area, making the statement true.
\item \textbf{Answer:} False
\\ \textit{Options:} A) True; B) False
\\ \textit{Note:} The statement is false because a compound is made up of two or more atoms that are chemically bonded together, whereas an atom is the smallest unit of an element and cannot exhibit the properties of a compound on its own.
\end{enumerate}

\section*{Long Form Response}
\begin{enumerate}[label=\arabic*.]
\setcounter{enumi}{10}
\item \textbf{Answer:} Volcanic eruptions in Earth's early history played a crucial role in shaping the planet's structure by releasing various materials such as gases, ash, and molten rock. These eruptions contributed to the formation of the Earth's crust and atmosphere, as gases like water vapor, carbon dioxide, and sulfur dioxide were expelled, helping to create a primordial atmosphere. The ash and lava from these eruptions solidified to form new landmasses and volcanic islands. This process is significant because it not only influenced the geological features of the planet but also set the stage for the development of ecosystems by creating diverse environments. Ultimately, these early volcanic activities were essential in establishing the foundational elements necessary for life to thrive on Earth.
\\ \textit{Note:} correct because it effectively outlines how volcanic eruptions released essential materials that contributed to the formation of the Earth's crust and atmosphere, highlighting the significance of these processes in shaping the planet's structure.
\item \textbf{Answer:} The speed of sound varies in different materials due to their physical properties, such as density and elasticity. Sound travels fastest in solids because their closely packed molecules allow sound waves to transmit energy quickly. For instance, sound travels faster in steel than in air because steel is denser and more elastic, enabling quicker vibrations. In contrast, gases, like air, have widely spaced molecules, which slows down the transmission of sound. Therefore, steel would allow sound to travel the fastest because its properties facilitate rapid energy transfer compared to other materials.
\\ \textit{Note:} correct because it accurately identifies that the speed of sound is influenced by the density and elasticity of materials, with sound traveling fastest in solids like steel due to their closely packed molecules facilitating quicker energy transmission compared to gases like air.
\end{enumerate}

\vfill
\noindent\textit{End of Answer Key}
\end{document}